\section{Demonstration}
\label{sec:demo}

\begin{figure}[t]
  \centering
  \placeholderfig{40mm}{figures/ui-main.png --- the main interface: the
  uploaded image with a selected region at left, the per-layer bias controls
  in the centre, and the unmodified and intervened predictions on the map at
  right}
  \caption{The GeoProbe interface. The map displays the unmodified
  prediction and the intervened prediction together with the geodesic
  distance separating them.}
  \label{fig:ui_main}
  \vspace{-3mm}
\end{figure}

A demonstration video of under three minutes is available at
\TODO{video URL}.

\subsection{Interaction}

Attendees operate the system directly, using either the images provided or
their own photographs, through three scenarios of increasing analytical
depth.

\subsubsection{Scenario 1: dependence upon an individual cue.}
The attendee uploads a street-level photograph and suppresses a single
region --- a shopfront sign or a licence plate --- by assigning a negative
in-region bias across all layers. The map is updated with the intervened
prediction and the resulting displacement. A displacement of several
thousand kilometres identifies a cue upon which the model demonstrably
depends; one of a few kilometres indicates that it does not, irrespective of
the prominence assigned to that region by saliency methods.
\TODO{specify the image used and the observed displacement}

\subsubsection{Scenario 2: comparison across cue categories.}
Rather than delineating a region manually, the attendee selects a concept
from the Grounding DINO vocabulary; the corresponding detections are loaded
as the region and the same suppression is applied. Iterating over
\texttt{street sign}, \texttt{vegetation}, \texttt{architecture} and
\texttt{license plate} for one image yields an ordering of the named cue
categories by their causal contribution to that prediction.

\subsubsection{Scenario 3: localisation of the effect within the network.}
The per-layer controls restrict the intervention to the early ($0$--$7$),
middle ($8$--$15$) or late ($16$--$23$) layer groups. Attendees observe that
the groups are not interchangeable
\TODO{state which group dominates, per Table~\ref{tab:groups}},
substantiating the proposition that geographic evidence is aggregated at a
particular depth rather than uniformly. The intervention is reversible
throughout: resetting the biases to zero restores the baseline prediction
exactly.

\subsection{Setup and Requirements}

GeoProbe operates entirely on a single laptop and requires no network
connectivity: the GeoCLIP weights, the $10^{5}$-point GPS gallery, the
precomputed location embeddings and the cached Grounding DINO detections are
all held locally, and a cached offline tile set substitutes for online map
tiles. We provide the machine (\TODO{GPU and memory}) and require only table
space, mains power and, where available, an external monitor. No specialised
equipment or network access is necessary.

\subsection{Relevance to the MMM Community}

The demonstration addresses three concerns central to this community:
multimodal vision--language models, whose internal behaviour increasingly
determines whether a deployed system may be trusted; explainability of a
causal rather than correlational character, which multimedia forensics
requires if attribution claims are to withstand scrutiny; and image
geolocation itself, whose dual use --- supporting the verification of
imagery on the one hand and the erosion of location privacy on the other ---
renders an understanding of the evidence these models exploit consequential
beyond the question of accuracy alone.
